\section{Параллельная реализация}

Все четыре программы написаны на C++17 и используют общий численный код.
Поэтому коэффициенты, начальное приближение и условие остановки во всех версиях совпадают.
Изменяется способ выполнения операций над локальными узлами и объединения результатов.

\subsection{OpenMP}

В версии OpenMP узлы одной прямоугольной сетки распределяются между потоками одного процесса.
Потоки работают с общей памятью и используют одни и те же массивы данных.
Двойные циклы для действия оператора, обновления векторов и вычисления скалярных произведений используют \code{parallel for collapse(2)} со статическим распределением итераций.
Суммы и максимумы вычисляются через \code{reduction}~\cite{omp}.
Каждый поток обновляет собственный набор элементов, поэтому одновременная запись в один узел не требуется.
Порядок сложения в редукции может отличаться от последовательного; побитовое совпадение результатов не предполагается.

\subsection{MPI}

MPI (Message Passing Interface) обеспечивает обмен сообщениями между процессами с раздельной памятью.
Каждый процесс хранит свою часть сетки.
MPI-процессы образуют прямоугольную решётку $P_x\times P_y$, где $P_xP_y=P$ и $P$ есть общее число процессов.
При выборе её размеров перебираются делители $P$.
Числа интервалов каждой координаты распределяются с остатком, поэтому размеры геометрических подобластей различаются не более чем на один интервал.
Среди допустимых разбиений выбирается ближайшее к квадрату по числу узлов.

Ограничение задания на отношение сторон применяется к узлам геометрической подобласти, включая общие интерфейсы.
Если в ней $m_x$ и $m_y$ интервалов, проверяется условие
\begin{equation}
    \frac12\leq\frac{m_x+1}{m_y+1}\leq2.
\end{equation}
При этом каждое внутреннее неизвестное имеет ровно одного владельца.
Например, при разбиении сетки $80\times80$ между двумя процессами геометрические подобласти имеют по $41\times81$ узлу с учётом общей границы.
Число принадлежащих процессу внутренних строк равно $39$ или $40$, поскольку всего таких строк $79$.

Для действия пятиузлового шаблона необходимы значения соседей за границей локального массива.
Они хранятся в дополнительном слое и обновляются перед каждым умножением на $\mat A$.
Обмен реализован четырьмя вызовами \code{MPI\_Sendrecv}; столбцы передаются с помощью производного типа данных~\cite{mpi}.
На обычной итерации выполняются два глобальных суммирования через \code{MPI\_Allreduce}: одно для шага $\alpha_k$, второе одновременно для нескольких скалярных величин после обновления решения.
Физические граничные узлы поправки остаются нулевыми.

\subsection{MPI+OpenMP}

Гибридная версия сохраняет MPI-разбиение и распараллеливает локальные циклы каждого процесса средствами OpenMP.
Число потоков задаётся одинаковым для всех процессов.
Все вызовы MPI выполняются главным потоком после окончания соответствующей параллельной области, что соответствует уровню \code{MPI\_THREAD\_FUNNELED}~\cite{mpi}.
Во время MPI-обмена остальные потоки процесса не выполняют локальные циклы решателя.
